\documentclass{article}
\usepackage{graphicx}
\usepackage{amsmath}
\usepackage{amsfonts}
\usepackage{bookmark}
\usepackage{booktabs}
\usepackage{exercises}
\graphicspath{ {./images/quiz3} }
\usepackage{indentfirst}
\usepackage{amsthm}
\usepackage{yhmath}
\usepackage{placeins}
\newcommand*{\oldneg}{\mathord{\sim}}
\newcommand\showdiv[1]{\overline{\smash{)}#1}}
\newcommand\Mydiv[1]{%
\smash{\raise.3ex\hbox{$\big)$}}\mkern-8mu
        \overline{\enspace\strut#1}}
\setlength{\tabcolsep}{5pt}
\theoremstyle{definition}
\newtheorem{definition}{Definition}

\setlength{\parindent}{.5in}
\setlength{\parskip}{12pt}
\renewcommand{\arraystretch}{2}
\newcommand{\?}{\stackrel{?}{=}}

\title{Test Answers (pg. 221)}
\begin{document}
\maketitle

\section*{1}
\[ -125^\frac{1}{3}=-(5^3)^\frac{1}{3}=-5 \]

\section*{2}
\[ 32^\frac{1}{5}=(2^5)^\frac{1}{5}=2 \]

\section*{3}
\[ \sqrt[4]{81}=\sqrt[4]{3^4}=3 \]

\section*{4}
\[ \sqrt[3]{27}=3 \]

\section*{5}
\[ 8^\frac{5}{3}=(2^3)^\frac{5}{3}=2^5=32 \]

\section*{6}
\[ 16^{-\frac{3}{2}}=\frac{1}{16^\frac{3}{2}}=\frac{1}{(4^2)^\frac{3}{2}}=\frac{1}{64} \]

\section*{7}
\[ (\sqrt[3]{-27})^2=-3^2=9 \]

\section*{8}
\[ (\sqrt[3]{64})^{-4}=(4)^{-4}=\frac{1}{4^4}=\frac{1}{256} \] 

\section*{9}
\[ \sqrt[3]{88}=\sqrt[3]{2^3\times 11}=2\sqrt[3]{11} \]

\section*{10}
\[ \sqrt[5]{16}\times \sqrt[5]{8}=\sqrt[5]{128}=\sqrt[5]{2^5\times 2^2}=2\sqrt[5]{4} \]

\section*{11}
\[ \sqrt{\frac{12}{49}}=\frac{\sqrt{12}}{\sqrt{49}}=\frac{2\sqrt{3}}{7} \]

\section*{12}
\[ \frac{\sqrt[3]{24}}{\sqrt[3]{9}}=\frac{2\sqrt[3]{3}}{\sqrt[3]{3^2}}\times\frac{\sqrt[3]{3}}{\sqrt[3]{3}}=\frac{2\sqrt[3]{9}}{3} \]

\section*{13}
\[ \sqrt[3]{64x^4y^2}=4x\sqrt[3]{y^2x} \]

\section*{14}
\[ \sqrt[4]{2x^6y^8z}=y^2x\sqrt[4]{2x^2z} \]

\section*{15}
\[ \sqrt[5]{\frac{x^6}{y^4}}=\frac{\sqrt[5]{x^6}}{\sqrt[5]{y^4}}=\frac{x\sqrt[5]{x}}{\sqrt[5]{y^4}}\times\frac{\sqrt[5]{y}}{\sqrt[5]{y}} \]
\[ \frac{x\sqrt[5]{xy}}{y} \]

\section*{16}
\[ \sqrt{\frac{75x^5y^6}{36xz^5}}=\sqrt{\frac{25x^4y^6}{12z^5}}=\frac{\sqrt{25x^4y^6}}{\sqrt{12z^5}} \]
\[ =\frac{5x^2y^3}{2z^2\sqrt{3z}}\times\frac{\sqrt{3z}}{\sqrt{3z}} \] \[ =\frac{5x^2y^3\sqrt{3z}}{6z^3} \]

\section*{17}
\[ f(x)+g(x)=2x+9+3x-1=5x+8 \]
The domain of this function is all real numbers, as the domains of $g$ and $f$ are also all real number.s

\section*{18}
\[ f(x)-g(x)=2x+9-(3x-1)=-x+10 \]
Again, the domain of this function is all real numbers.

\section*{19}
\[ f(x)\times g(x)=(2x+9)(3x-1)=6x^2+25x-9 \]
Again, the domain of this function is all real numbers.

\section*{20}
\[ \frac{f(x)}{g(x)}=\frac{2x+9}{3x-1} \]
The domain of this function is all real numbers other than $\frac{1}{3}$, as $x=\frac{1}{3}$ would make $g(x)=0$, and we cannot divide by zero.

\section*{21}
\[ f(g(x))=f(3x-1)=2(3x-1)+9=6x-2+9=6x+7 \]
Again, the domain of this function is all real numbers.

\section*{22}
\[ g(f(x))=g(2x+9)=3(2x+9)-1=6x+27-1=6x+26 \]
Again, the domain of this function is all real numbers.

\section*{23}
\[ f(f(x))=f(2x+9)=2(2x+9)+9=4x+18+9=4x+27 \]
Again, the domain of this function is all real numbers.

\section*{24}
\[ g(g(x))=g(3x-1)=3(3x-1)-1=9x-3-1=9x-4 \]
Again, the domain of this function is all real numbers.

\section*{25}
\[ y=-2x+5 \] \[ x=-2y+5 \textrm{ (Switch variables)} \] \[ x-5=-2y \] \[ y=-\frac{x-5}{2} \]

\section*{26}
\[ y=\frac{1}{3}x+4 \] \[ x=\frac{1}{3}y+4 \textrm{ (Switch variables)} \] \[ x-4=\frac{1}{3}y \] \[ y=3x-12 \]

\section*{27}
\[ y=5x-12 \] \[ x=5y-12 \textrm{ (Switch variables)} \] \[ y=\frac{x+12}{5} \]

\section*{28}
\[ y=\frac{1}{2}x^4,x\ge 0 \] \[ x=\frac{1}{2}y^4,y\ge 0 \textrm{ (Switch variables)} \]
\[ 2x=y^4,y\ge 0 \] \[ \pm\sqrt[4]{2x}=y \] Because $y\ge 0$, the correct inverse function is $+\sqrt[4]{2x}=y$.

\section*{29}
\[ y=x^3+5 \] \[ x=y^3+5 \] \[ x-5=y^3 \] \[ y=\sqrt[3]{x-5} \]

\section*{30}
\[ y=-2x^3+1 \] \[ x=-2y^3+1 \] \[ -\frac{x-1}{2}=y^3 \] \[ y=\sqrt[3]{-\frac{x-1}{2}} \]

\section*{31}

\section*{34}
\[ \sqrt{3x+7}=4 \] \[ 3x+7=16 \] \[ x=3 \]
Check for extranoues solutions:
\[ \sqrt{3(3)+7}\? 4 \] \[ \sqrt{16}\? 4 \] \[ 4=4 \]

\section*{35}
\[ \sqrt{3x}-\sqrt{x+6}=0 \]
\[ \sqrt{3x}=\sqrt{x+6} \]
\[ 3x=x+6 \] \[ 2x-6=0 \] \[ x=3 \]
Check for extraneous solutions:
\[ \sqrt{3(3)}-\sqrt{3+6}\? 0 \] \[ 9-9=0 \] 

\section*{36}
\[ x-3=\sqrt{x-1} \] \[ x^2-6x+9=x-1 \] \[ x^2-7x+10=0 \] \[ (x-5)(x-2)=0 \] \[ x=5\vee x=2 \]
Check for extraneous solutions:
\[ 5-3\? \sqrt{5-1} \] \[ 2=2 \] and \[ 2-3\? \sqrt{2-1} \] \[ -1\ne 1 \] thus meaning $x=5$ is the only solution.

\section*{37}
(a): \[ E=625s^2 \] \[ \frac{E}{625}=s^2 \] \[ s=\frac{\sqrt{E}}{25} \]
(b): \[ s=\frac{\sqrt{120000}}{25}\approx 13.9 \] meters per second. Convert to kilometers per hour:
\[ \frac{13.9\; meters}{1\; second}\times \frac{3600\; seconds}{1\; hour}=\frac{50040\; meters}{1\; hour} \]
\[ =\frac{50040\; meters}{1\; hour}\times\frac{0.001\; kilometers}{1\; meters}=\frac{50.004\; kilometers}{1\; hour} \]
(c): Suppose we have a given value $E$ for the kinetic energy. Thus \[ s=\frac{\sqrt{E}}{25} \] Suppose we then take $2E$. Thus:
\[ s=\frac{\sqrt{2E}}{25} \] This means that speed will not double, as $\sqrt{2E}<2\sqrt{E}$ (as $2\sqrt{E}=\sqrt{4E}$).

\section*{38}
(a): \[ h=0.9(200-a) \] \[ \frac{10h}{9}=200-a \] \[ a=200-\frac{10h}{9} \]
(b): Suppose $h=36$. Thus: \[ a=200-\frac{10(36)}{9}=200-40=160 \]

\end{document}

 
